\section*{Chapitre \ref{ch:b2:himalaya-andes} --- Chaînes de collision et chaînes andines comparées}

\begin{solution}{exo:b2:himalaya-andes:1}
(a) $r = h\rho_c/(\rho_m - \rho_c) = 3.0\times2700/600 = \qty{13.5}{km}$. (b) $T = 35 + 3.0 +
13.5 = \qty{51.5}{km}$, soit $35 + 5.5\times3.0$. (c) $h = \Delta T(\rho_m - \rho_c)/\rho_m =
10\times600/3300 = \qty{1.8}{km}$.
\end{solution}

\begin{solution}{exo:b2:himalaya-andes:2}
(a) Avant-fosse indo-gangétique~; chevauchement frontal principal~; sous-Himalaya (Siwaliks)~; chevauchement de
bordure principal~; Petit Himalaya~; chevauchement central principal~; séquence du Grand Himalaya~; système
de détachement du Tibet méridional~; séquence téthysienne de l'Himalaya~; puis la suture de l'Indus--Tsangpo
et la croûte asiatique. (b) Le système de détachement du Tibet méridional. (c) Le degré de \omterm{def:b2:metamorphic-facies:metamorphism}{métamorphisme} augmente
vers le nord~: \omterm{def:b2:alpine-orogen:foreland}{molasse} non métamorphique, Petit Himalaya de \omterm{def:b2:metamorphic-facies:metamorphism}{métamorphisme} moyen, séquence du Grand Himalaya
de haut degré et en partie fondue~; au-dessus du STDS, il chute brusquement vers les sédiments téthysiens
non métamorphiques ou faiblement métamorphiques, puisque la faille normale a retiré les roches intermédiaires.
\end{solution}

\begin{solution}{exo:b2:himalaya-andes:3}
(a) $196.5\ \text{m}/9200\ \text{an} = \qty{0.02136}{m}$ par an, soit environ \qty{21.4}{mm}
par an. (b) C'est \qty{20}{\percent} de plus que les \qty{17.8}{mm} par an du
GPS, mais du même ordre~: un taux est moyenné sur neuf mille ans, l'autre sur vingt
ans. (c) Presque tout le raccourcissement à travers la chaîne est absorbé au front~: le
MBT, le MCT et les autres failles situées au nord du MFT sont, à cette longitude,
pratiquement inactifs~; elles s'enracinent toutes dans le même MHT, qui affleure au MFT.
\end{solution}

\begin{solution}{exo:b2:himalaya-andes:4}
(a) Collision~: la suture marque l'océan disparu. (b) Type andin~: un arc a besoin
d'une plaque qui plonge. (c) Collision, en général~: fusions crustales d'une croûte
épaissie et chauffée (il en existe aussi dans la croûte andine épaissie, mais les \omterm{def:b2:magma-emplacement:magma}{magmas}
d'arc y dominent). (d) Type andin. (e) Les deux~: les Siwaliks et l'avant-pays subandin.
(f) Les deux~: le MHT dans l'Himalaya, l'interface de subduction dans les Andes.
\end{solution}

\begin{solution}{exo:b2:himalaya-andes:5}
(a) $17.8/45 = \qty{40}{\percent}$ à l'est, $20.5/45 = \qty{46}{\percent}$ à l'ouest~:
environ la moitié. (b) Dans l'épaississement et les failles du Tibet et du reste de l'Asie,
et dans l'extrusion vers l'est de blocs le long des grands décrochements. (c) $45 - 17.8 =
\qty{27.2}{mm}$ par an, soit \qty{272}{km} en \qty{10}{Ma}~;
$45 - 20.5 = \qty{24.5}{mm}$ par an, soit \qty{245}{km}.
\end{solution}

\begin{solution}{exo:b2:himalaya-andes:6}
(a) $r = 4.5\times3.7 = \qty{16.65}{km}$~;
$\Delta F = 2700\times9.8\times3700\times(35\,000 +
1850 + 8325) = \qty{4.42e12}{N.m^{-1}}$. (b) Pour le Tibet~: \qty{6.12e12}{N.m^{-1}}~; l'Altiplano en a
\qty{72}{\percent}. (c) Avec $r = 4.5h$, $\Delta F = \rho_c g\,(Hh + 2.75h^2)$~: pour $h \ll H$, le terme
$Hh$ domine et $\Delta F$ est proportionnel à $h$~; pour un plateau élevé, le terme en $h^2$ croît plus
vite. Chaque kilomètre d'altitude supplémentaire coûte de plus en plus de poussée aux bords, jusqu'à
ce que les forces disponibles ne puissent plus élever le plateau~; il s'étale alors au lieu de s'élever.
\end{solution}

\begin{solution}{exo:b2:himalaya-andes:7}
(a) $2700\times9.8\times8848.86/2 = \qty{1.17e8}{Pa} = \qty{117}{MPa}$. (b) À
\qty{100}{MPa}~: $0.85\times100 = \qty{85}{MPa}$~; à \qty{200}{MPa}~: $0.85\times200 = 50 +
0.6\times200 = \qty{170}{MPa}$, les deux branches s'accordent. (c) $h = 2\sigma_d/(\rho g) =
2\times1.70\times10^8/(2700\times9.8) = \qty{12.8}{km}$~: un ordre de
grandeur, de l'ordre des sommets les plus hauts, pas loin au-dessus d'eux.
\end{solution}

\begin{solution}{exo:b2:himalaya-andes:8}
(a) $11\times111.19\times\cos\ang{30} = \qty{1059}{km}$.
(b) $9\times10^{-3}/1.059\times10^{6}
= \num{8.5e-9}$ par an (entre \num{6.6e-9} et \num{10.4e-9} pour 7 et \qty{11}{mm} par an), soit
$8.5\times10^{-9}/3.156\times10^{7} = \qty{2.7e-16}{s^{-1}}$. (c) Les deux intervalles se recouvrent~: $(12.4 \pm 4.7) \times 10^{-9}$ couvre de 7.7 à
\num{17.1e-9}. Le plateau méridional au nord du Népal s'étire aussi vite que la chaîne dans son ensemble, ou un peu plus vite.
\end{solution}

\begin{solution}{exo:b2:himalaya-andes:9}
(a) $0.9/6.6 = \qty{14}{\percent}$. (b) $M_0 = 10^{1.5\times8.5 + 9.1} = 10^{21.85} =
\qty{7.1e21}{N.m}$. (c) $7.1\times10^{21}/6.6\times10^{19} = 107$~: un tel séisme environ
tous les 110 ans sous le Népal, bien plus que ne l'enregistre l'histoire~; soit de bien
plus grands séismes sont à attendre, soit une partie du déficit se libère autrement.
\end{solution}

\begin{solution}{exo:b2:himalaya-andes:10}
(a) Prendre la profondeur de compensation à la
base de la croûte de l'Altiplano, $61.5 - 3.7 =
\qty{57.8}{km}$ sous le niveau de la mer. La croûte de la Puna atteint $50.5 - 4.4 = \qty{46.1}{km}$, si bien
que la colonne de la Puna contient \qty{11.7}{km} de manteau au-dessus de cette profondeur~: $2800\times61.5 =
\rho_P\times50.5 + 3300\times11.7$. (b) $\rho_P = (172\,200 - 38\,610)/50.5 =
\qty{2645}{kg.m^{-3}}$. (c) La différence, \qty{155}{kg.m^{-3}}, se place à l'extrémité basse de la fourchette de 150 à
\qty{200}{kg.m^{-3}} citée~: la croûte seule aurait besoin d'une croûte de la Puna anormalement légère. Une colonne plus légère dessous y
parvient plus facilement~: une lithosphère mantellique retirée par délamination, remplacée par une asthénosphère chaude et légère.
\end{solution}

\begin{solution}{exo:b2:himalaya-andes:11}
(a) $g = 105/\tan\ang{30} = \qty{181.9}{km}$. (b) $g = 90/\tan\ang{30} + 15/\tan\ang{4} =
155.9 + 214.5 = \qty{370.4}{km}$~: l'arc se déplace de \qty{188.5}{km} vers l'intérieur. (c) La plaque n'atteint
jamais \qty{105}{km} sous la plaque chevauchante sur ces \qty{300}{km}, et aucun coin de manteau chaud ne se trouve
au-dessus d'elle~: aucun \omterm{def:b2:magma-emplacement:magma}{magma} n'est produit, et l'arc s'éteint jusqu'à ce que la plaque replonge, plus à l'est.
\end{solution}

\begin{solution}{exo:b2:himalaya-andes:12}
(a) \qty{21}{mm} par an font \qty{21}{km} par million d'années~: $21\times37.5 =
\qty{787.5}{km^3}$ par kilomètre de chaîne et par million d'années. (b) $300/787.5 =
\qty{38}{\percent}$. (c) Dans l'épaississement de la chaîne et du sud du Tibet, dans la croûte
entraînée sous le Tibet avec l'Inde, et dans la matière que le MHT fait passer vers le nord.
\end{solution}

\begin{solution}{pb:b2:himalaya-andes:1}
\textbf{1.} Sept échantillons de \ang{28.0} à \ang{37.0}~N~: $34\,550/7 = \qty{4936}{m}$, $h_T
= \qty{4.9}{km}$.
\textbf{2.} Cinq échantillons de \ang{-69} à \ang{-65}~: $18\,650/5 = \qty{3730}{m}$, $h_A =
\qty{3.7}{km}$, la valeur citée pour l'Altiplano~; la Puna, plus au sud, est plus haute.
\textbf{3.} Tibet~: au moins $9\times111.19 = \qty{1001}{km}$~; l'Altiplano~: au moins $4\times
104.5 = \qty{418}{km}$.
\textbf{4.} $(5000 - 70)/(1.5\times111\,190) = 0.0296$~: \qty{3.0}{\percent}.
\textbf{5.} Le profil andin commence sur le fond océanique et dans la fosse, à plusieurs
kilomètres sous le niveau de la mer, à l'ouest du plateau~; le profil tibétain commence sur
une plaine continentale près du niveau de la mer, et se termine sur le haut bassin du Tarim.
\textbf{6.} $r = 4.5\times4.9 = \qty{22.05}{km}$~; $T_T = 35 + 4.9 + 22.05 = \qty{62.0}{km}$.
\textbf{7.} $r = 4.5\times3.7 = \qty{16.65}{km}$~; $T_A = 35 + 3.7 + 16.65 = \qty{55.4}{km}$.
\textbf{8.} Tibet~: \qty{62.0}{km} contre 70 à \qty{75}{km}, soit \qty{10.5}{km} de moins que la valeur
médiane \qty{72.5}{km} (\qty{14}{\percent})~; l'Altiplano~: \qty{55.4}{km} contre 59 à \qty{64}{km},
soit \qty{6.1}{km} de moins que \qty{61.5}{km} (\qty{10}{\percent}). Le Tibet est le plus éloigné.
\textbf{9.} $72.5 - 35 = 4.9\times3300/(3300 - \rho_c)$~: $3300 - \rho_c = 16\,170/37.5 =
431.2$, $\rho_c = \qty{2869}{kg.m^{-3}}$.
\textbf{10.} $26.5 = 3.7\times3300/(3300 - \rho_c)$~: $3300 - \rho_c = 460.8$, $\rho_c =
\qty{2839}{kg.m^{-3}}$.
\textbf{11.} Une partie de la charge est portée par la flexion de la plaque, et non par une racine locale~; la
croûte normale n'avait peut-être pas \qty{35}{km} d'épaisseur~; le manteau sous-jacent peut être plus léger ou
plus dense que supposé~; le Moho sismique peut ne pas être la base de la croûte au sens de la masse volumique.
\textbf{12.} La base de la croûte tibétaine se trouve à $72.5 - 4.9 = \qty{67.6}{km}$, celle de l'Inde
à \qty{37.5}{km}~: $t = \qty{30.1}{km}$. $\Delta g = 2\pi\times6.674\times10^{-11}\times370\times
30\,100 = \qty{4.67e-3}{m.s^{-2}}$, soit \qty{-467}{mGal}
par rapport à l'Inde, proche des \qty{460}{mGal} observés.
\textbf{13.} $W_0 = 440\times60/35 = \qty{754.3}{km}$~; $S = \qty{314.3}{km}$.
\textbf{14.} $(314.3 - 250)/314.3 = \qty{20}{\percent}$,
le déficit constaté dans les Andes centrales.
\textbf{15.} $0.55\times314.3 = \qty{173}{km}$, moins que les
\qty{250}{km} des \omterm{def:b2:alpine-orogen:balanced}{coupes équilibrées}, qui suffiraient alors largement.
\textbf{16.} La couverture raccourcie au-dessus du \omterm{def:b2:alpine-orogen:decollement}{décollement} doit correspondre à un raccourcissement
du socle situé dessous, plus à l'ouest, puisque les deux doivent perdre la même longueur à travers la
chaîne~; le raccourcissement du socle épaissit la croûte sous la Cordillère orientale et le plateau.
\textbf{17.} $314.3/30 = \qty{10.5}{mm}$ par an,
$10.5/72 = \qty{15}{\percent}$ de la convergence.
\textbf{18.} Elle est consommée par la subduction~: la plaque de Nazca plonge sous le continent.
\textbf{19.} $r = \qty{22.05}{km}$~; $\Delta F = 2700\times9.8\times4900\times(35\,000 + 2450
+ 11\,025) = \qty{6.28e12}{N.m^{-1}}$.
\textbf{20.} $2700\times9.8\times3700\times45\,175 = \qty{4.42e12}{N.m^{-1}}$~;
le rapport vaut 0.70.
\textbf{21.} $6.28\times10^{12}/10^{5} = \qty{6.3e7}{Pa}$, soit environ \qty{63}{MPa}.
\textbf{22.} $9\times10^{-3}/1.06\times10^{6} = \num{8.5e-9}$ par an.
\textbf{23.} La poussée de l'Inde agit du nord au sud et, sur le haut plateau, est équilibrée par le
propre poids du plateau~: dans la direction nord--sud, le plateau ne peut plus s'épaissir, si bien que
la contrainte verticale, le poids, est la plus grande~; la contrainte la plus faible est est--ouest,
vers les bords libres de l'Asie, et le plateau s'étire dans cette direction, sur des rifts nord--sud.
\textbf{24.} \textbf{\qty[text-series-to-math=true]{62}{km}},
soit environ \qty{86}{\percent} des \qty{72.5}{km} mesurés.
\end{solution}
